\renewcommand{\theequation}{E\thechapter.\arabic{equation}}
\label{sec:appendixE}

The spatially resolved model is potential future work that could extend from the 0D flow model and analyses in this work.
The 0D flow model will rapidly generate improved \gls{DNP} group parameters, but it simplifies the treatment of flow.
Instead of modeling the actual fuel salt flow, it will estimate fuel circulation using residence times for each region.
A spatially resolved model would be useful for capturing the effects that the 0D flow model only capture with residence times for specific \gls{MSR} geometries.
Figure \ref{fig:sophisticated-model} shows the flow diagram for the spatially resolved model, which could model irradiation of a target fissile sample using Moltres.

\begin{figure}[H]
    \centering
    \includegraphics[width=0.9\linewidth]{images/spatial-model.png}
    \caption{Flow diagram of spatially resolved flow model used to generate improved \gls{DNP} group parameters.}
    \label{fig:sophisticated-model}
\end{figure}

The spatially resolved flow model introduces two new parameters absent in the 0D flow model.
The first is the mesh, which defines the flow loop geometry.
MOOSE provides built-in meshing tools, though complex geometries may require external tools such as Gmsh \cite{geuzaine2009gmsh}.
The second is the thermal-hydraulics parameters, which capture physical phenomena enabled by spatial resolution, as discussed in Section \ref{sec:moltres}.
Additionally, since the depletion would happen in Moltres instead of OpenMC, OpenMC does not need as much data.

In the spatially resolved model, the in-core and ex-core residence times would be outputs rather than inputs.
These times derive from the movement of the irradiated sample over time in the fuel salt.
The in-core and ex-core residence times vary as the irradiated sample diffuses and mixes further into the surrounding salt.
The net residence time, $\tau_{net}$, remains constant unless the model changes, such as an increased temperature difference or induced flow rate.
Equation \eqref{eq:incore-time} defines the in-core residence time as:
\begin{equation}
    \tau_{in} (t) = \frac{\int_{V_{in}} \int_{0}^t N(\vec{x}, \tau) d\tau d\vec{x}}{\int_{V} \int_{0}^t N(\vec{x}, \tau) d\tau d\vec{x}} \tau_{net},
    \label{eq:incore-time}
\end{equation}
where
\begin{align}
\tau_{in} &= \text{ the in-core residence time $[s]$} \nonumber \\
V_{in} &= \mbox{ the in-core volume $[cm^3]$} \nonumber \\
V &= \mbox{ the total primary volume $[cm^3]$} \nonumber \\
N(\vec{x}, t) &= \mbox{ the time and space resolved sample concentration $[\# \cdot cm^{-3}]$} \nonumber \\
\tau_{net} &= \mbox{ the net residence, or circulation, time of the fuel salt $[s]$.} \nonumber
\end{align}

The ex-core residence time can be calculated in the same way, or alternatively as shown in Equation \eqref{eq:excore-time} as:
\begin{equation}
    \tau_{ex} = \tau_{net} - \tau_{in}.
    \label{eq:excore-time}
\end{equation}